% ============================================================================
%  02_method.tex --- Sec. 2, Energy-Value Supervision.  No floats.
%
%  2026-08-30 REFRAME.  Sec. 2 is the METHOD section of a method paper.
%  2026-08-30 RELATED-WORK PASS.  A compact related-work subsection was
%  PROMOTED from Appendix C (sections/A3_related.tex) and is now 2.1, so the
%  positioning is in the main text where a reviewer will find it.  Sec. 2 now
%  has five subsections:
%    2.1  related work: three groups + the intersection sentence   <-- promoted
%    2.2  the bond-free flow and its clamped positional density
%    2.3  the group-wise value objective          <-- BOXED, visually dominant
%    2.4  why the partition function cancels      <-- and why grouping is forced
%    2.5  a gradient-supervision baseline         <-- comparator, NOT a component
%
%  The proposed objective is the only boxed display in the paper.  The force
%  baseline in 2.5 must never be boxed or otherwise made to compete with it.
%
%  ATTRIBUTION IS MADE IN 2.1 AND NOWHERE ELSE.  The log-variance divergence,
%  the cancellation of Z(c) and off-policy estimation are prior work; 2.4
%  states the technical consequence and cross-references 2.1 for the credit.
%  Do NOT restore the citation list in 2.4 -- that would double-credit.
%
%  NOTATION.  lambda_E and L_value throughout; the historical code field is
%  named once, in 2.2.  kT_eff, never T_eff.  log q_theta, never log p_theta.
%
%  2026-08-19.  The eps = 1/(2n) SMOOTHED-MARGINAL READING IS WRONG AND HAS
%  BEEN DELETED (paper/ESTIMATOR_SEMANTICS.md).  Do not reinstate it.
% ============================================================================

\section{Energy-Value Supervision for De Novo Molecular Flows}
\label{sec:method}

\subsection{Related work: two literatures, and the question between them}
\label{sec:related}

\paragraph{De novo 3D molecular generation.} Equivariant diffusion and flow-matching models
generate positions and types jointly, from EDM and GeoLDM \citep{edm,geoldm} through
graph-and-geometry hybrids \citep{midi,symphony,semlaflow} to FlowMol3 \citep{flowmol3}, the
backbone used here, and Zatom-1 \citep{zatom1}, the one published generator trained on the same
corpus; EBMol \citep{ebmol} learns a scalar energy from data rather than from a physical potential.
These models are judged on validity, on sample quality, or on an energy the model itself learns, and
they define the state of the art at the task they set out to do. What such evaluations generally do
not ask is whether the tractable density the model defines responds to an \emph{external} physical
energy. That is a different question rather than a shortfall --- it is not the question these papers
pose --- and we claim superiority over none of them.

\paragraph{Conditional conformer and ensemble models.} Torsional diffusion \citep{torsionaldiff},
the GEOM conformer literature \citep{geom} and DECAF \citep{decaf} take the molecular graph, or the
molecular identity, as input and return conformers for it. The composition is supplied rather than
generated, so the density they define is conditional on a fixed chemical graph. That input is what
makes theirs a different problem from a de novo generator's, and why their protocols do not transfer
to one unchanged.

\paragraph{Energy-targeted samplers, and where our objective comes from.} Samplers that target
$e^{-E/kT}$ directly \citep{noe2019boltzmann,tbg,sbg,arbg,adjointsampling,idem,ewfm} obtain
equilibrium densities; mass-covering correction by annealed importance sampling \citep{fab}
addresses the cross-group failure that \Cref{cor:basins} formalises, potential-score matching
\citep{psm} supervises with the potential's gradient instead of its value, and \citet{idflow} drive
a predictor--refiner with an \emph{internal, learned} energy instead of an external one. Each takes
the system, the alphabet or the molecular graph as given: Adjoint Sampling \citep{adjointsampling},
the closest use of a universal potential as a training signal, trains against an eSEN energy but
samples conformers of a \emph{given} molecular graph. They therefore solve a stronger sampling
problem --- equilibrium, not ordering --- on a state space that is narrower and fixed in advance.
Our value term is inherited from this literature and is not claimed as new:
$\Var[\log q_\theta + E/kT]$ is the \emph{log-variance divergence} named by \citet{vargrad}, placed
in a path-space family by \citet{nusken2021pathspace}, given an optimal-control reading by
\citet{berner2024optimalcontrol} and established for diffusion samplers by
\citet{richter2024improvedsampling}; that a variance annihilates the normalising constant, and that
the divergence may be estimated off-policy under an arbitrary reference measure
\citep{sendera2024offpolicy}, are structural properties of that family, and we claim neither.

\paragraph{The intersection.} Existing de novo generators decide what molecule and what geometry to
produce, but do not test the tractable density they define against a transferable physical energy;
energy-targeted samplers test or optimise exactly that relation, but assume the molecular system or
the graph. This paper works at the intersection, asking what group-wise energy \emph{values} from a
universal potential do to the density of a generator that is still choosing its molecule. To our
knowledge the combination studied here --- an external universal potential supervising the
tractable density of a bond-free de novo flow --- has not been evaluated in this form.
\Cref{app:related} gives the formal setting and the fuller discussion.

\subsection{The bond-free flow and its clamped positional density}
\label{sec:method:flow}

A molecular state is a pair $(c,x)$: a composition $c$, the element multiset together with the total
charge, and a centre-of-mass-free conformation $x \in \R^{3N}$. The backbone is FlowMol3
\citep{flowmol3}, an SE(3)-equivariant conditional flow-matching model
\citep{lipman2023flowmatching,liu2023rectifiedflow,albergo2023interpolants} with
$x_t = (1-t)x_0 + t x_1$, $x_0 \sim \mathcal{N}(0,\sigma^2 I)$ at $\sigma = 1$, and
$\Ls_{\mathrm{FM}}$ the velocity-regression loss; types and charges ride a continuous-time Markov
chain. \emph{De novo} describes that backbone and its task; every energy measurement here is instead
conditional on a fixed endpoint composition, types and charges held at their data values whenever
the teacher is queried. The teacher is the universal neural potential trained on OMol25
\citep{omol25,esen}: 83 elements, DFT energies and forces, no bond labels, making the generator
bond-free by necessity, bonds recovered post hoc where an evaluation needs them \citep{xyz2mol}. It
supplies $E(c,x)$ in eV and $F = -\nabla_x E$, so
$\pi(x \mid c) = e^{-E(c,x)/kT}/Z(c)$; neither $Z(c)$ nor $\sum_c Z(c)$ is ever formed, and
$kT = 1.0$~eV is a \emph{numerical} target that keeps $E/kT$ tame, not a physical temperature.
Queries pass charge only and default to a singlet, so every evaluation below is on main-group
organic parents (\Cref{ass:charge}).

Because the generator is a continuous normalising flow, its log-density follows from the
instantaneous change of variables \citep{ffjord}: integrating the field backwards from $x_1$ and
accumulating the divergence gives
\begin{equation}
  \log q_\theta(x_1 \mid c) \;=\; \log p_0(x_0) - \int_0^1 \nabla \!\cdot\, v_\theta(x_t,t)\,\mathrm{d}t ,
  \label{eq:ffjord}
\end{equation}
with the trace estimated by Rademacher--Hutchinson probes \citep{hutchinson}. Throughout,
$q_\theta(x \mid c)$ is the \emph{clamped positional density}: \eqref{eq:ffjord} integrated with the
discrete channels, atom types, formal charges and the empty bond tensor, held at their endpoint
labels for the whole reverse trajectory, only the positions evolving. It costs one reverse solve per
geometry, is normalised on the centre-of-mass-free subspace by construction, and is the object this
paper supervises and evaluates. \textbf{It need not equal the joint generator's
conditional density $p_\theta(x \mid c)$}: the two coincide only if the discrete trajectory is
determined by its endpoint, and we have not quantified the difference (\Cref{ass:A2}).

\subsection{The group-wise value objective}
\label{sec:method:obj}

The intervention constrains the \emph{value} of that log-density rather than its gradient. The
Boltzmann residual $w(x) = \log q_\theta(x \mid c) + E(c,x)/kT$ is constant exactly where
$q_\theta(\cdot \mid c) \propto e^{-E(c,\cdot)/kT}$, so asking only that $w$ be \emph{flat} across a
set of geometries asks for the physics and nothing else. We pre-compute shards of $K$ perturbations
per parent, each carrying its teacher energy; every step draws $M$ parents and penalises the
variance of the residual inside each group:
\begin{equation}
  \boxed{\;
  \begin{aligned}
    \Ls_{\mathrm{ours}} \;&=\; \Ls_{\mathrm{FM}} \;+\; \lambda_E \,\Ls_{\mathrm{value}} ,
    \qquad
    \Ls_{\mathrm{value}} \;=\; \frac{1}{M}\sum_{m=1}^{M} \Var_k\bigl( w_{m,k} \bigr) ,
    \\[1pt]
    w_{m,k} \;&=\; \log q_\theta(x_{m,k} \mid c_m) \;+\; E_{m,k}/kT .
  \end{aligned}
  \;}
  \label{eq:lvalue}
\end{equation}
where $w_{m,k}$ is the residual at the $k$-th geometry of parent $m$, evaluated with that parent's
own composition $c_m$ and pre-computed teacher energy $E_{m,k}$. We use
$\lambda_E = 3\times10^{-5}$ (code field: \texttt{lambda\_2}), zero over a flow-matching warm-up and
then ramped, applied every few steps because each evaluation is a reverse-time solve carrying second
derivatives (\Cref{app:hyper}). Nothing else changes: same backbone, optimiser and data, no bond
labels, and the model we report as ours sets $\lambda_1 = 0$ and the anchor weight
$\lambda_3 = 0$.

\subsection{Why the partition function cancels, and why grouping is mandatory}
\label{sec:method:cancel}

Three properties of \eqref{eq:lvalue} do the work. First, $\log Z(c)$ enters $w$ as an additive
constant and a variance annihilates additive constants, so the normaliser cancels within a group and
is never formed --- the reason values from a transferable potential are usable at all, since $Z(c)$
changes with every composition. Second, grouping is mandatory rather than convenient, precisely
because that constant differs across molecules: its spread is set by size and composition rather
than by deviation from the Boltzmann law, so a variance over a batch of \emph{different} molecules
would measure the spread of $\log Z$ instead (\Cref{rem:crossbatch}).
Third, and separately from the cancellation, \eqref{eq:lvalue} is a group-wise estimate, against the
unnormalised Gibbs weight, of the \emph{log-variance divergence}, which may be estimated under an
\emph{arbitrary} reference measure. That, and not the cancellation, is what makes the supervision
off-policy and pre-computed shards valid: no sample need come from the model. The divergence, the
cancellation of $Z(c)$ inside it, and its off-policy estimation are prior work, credited in
\Cref{sec:related}. Ours is the deployment: a group-conditioned use of that estimator inside a mixed
discrete--continuous, bond-free de novo generator, at the price of one free constant per group
(\Cref{sec:theory}).

\paragraph{Practicalities.}
\label{sec:method:impl}
\emph{(i)} $\lambda_E$ was fixed once from a smoke run before any Boltzmann evaluation and never
tuned against a reported metric; the weight is small, its \emph{contribution} is not.
\emph{(ii)} \eqref{eq:ffjord} is integrated in $n$ explicit Euler steps over the whole of $[0,1]$,
the divergence integral by the midpoint rule on the same interval; the interval is not truncated, so
$n$ controls discretisation error rather than which density is meant (\Cref{app:estimator}).
Training uses $n = 4$ and every reported evaluation $n = 12$. \emph{(iii)} The reverse-time solve back-propagates second derivatives through a stochastic
trace, and runs carrying it reach non-finite weights at a rate that runs without it do not ($6$ of
$20$ against $0$ of $10$ in the matched grid of \Cref{sec:ablation}); a guard aborts on the first
non-finite parameter, and because divergence is not independent of the objective, \textbf{every
value-cell mean we report is conditional on completion}, with launched, completed and diverged
counts beside it. \emph{(iv)} The gradient through $\Ls_{\mathrm{value}}$ is a
\emph{truncated-adjoint surrogate}: the reverse trajectory is advanced without gradient tracking,
each state detached from the previous, so gradient flows only through the divergence evaluations; we
do not differentiate through the solve and do not claim exact likelihood training
(\Cref{app:asimplemented}).

\subsection{A gradient-supervision baseline}
\label{sec:forcebaseline}

The obvious alternative to values is the potential's gradient, measured separately under the same
backbone and budget. Flow matching gives the score of its time-$t$ marginal in closed form
(\Cref{lem:score}), $s_\theta(x_t,t) = [\,t\, v_\theta(x_t,t) - x_t\,]/[(1-t)\,\sigma^2]$, which at
the Boltzmann optimum equals $F/kT$, so a natural comparator regresses one on the other,
\begin{equation}
  \Ls_{\mathrm{force}}
  = \E_{t^{*},(c,x_1)}\Bigl[\, N(c)^{-1} \textstyle\sum_{i}
  \ell\bigl( s_\theta(x_{t^{*}}, t^{*})_i - \tfrac{F(c,x_1)_i}{kT} \bigr) \Bigr] .
  \label{eq:score}
\end{equation}%
The read-out has a pole at $t = 1$, so we probe only at $t^{*} \in \{0.85, 0.92, 0.97\}$ and average
a per-atom Huberised, norm-capped error $\ell$ (\Cref{app:constants}). The force target is read at
the data endpoint $x_1$ while the score is read on the conditional path, so this is a naive
off-policy endpoint form of force matching, a comparator and not a component of \eqref{eq:lvalue}: the model we report as ours has
$\lambda_1 = 0$, and results from $\lambda_1 > 0$ are labelled as such (\Cref{sec:ablation}).
